% 07-conclusion.tex — 结论

\section{结论}
\label{sec:conclusion}

本文给出 FamilyMask 多族 CA 编码与 $maze\_quality$ 八维几何质量度量两条主线, 用 15-pattern 对照基准和两个尺度的演化扫描验证了两者的协同效果。

\noindent\textbf{FamilyMask 编码空间.}
FamilyMask 把经典 B/S 型规则推广为最多 16 族并行、按优先级仲裁的离散结构, 搜索空间从 $2^{18}$ 跃升至 $2^{1648}$ (16 slot $\times$ 103 bits/slot, 103 = 1 active + 80 cells + 9 birth + 9 survive + 4 priority)。单族搜索空间 $2^{103}$, 全 16 族染色体 $2^{1648}$。当某个随机种子恰好落在 ``对的多族组合'' 邻域时, 多族可突破单族天花板 --- maxFam=8 的峰值 0.8233 高于 maxFam=1 的 0.7999 即此机制。

\noindent\textbf{maze\_quality 度量.}
$maze\_quality$ 用 8 维连续可微子度量 + 两层 $\min$-聚合 + 墙比软门控, 在 15-pattern 基准的 9 个伪迷宫反例上全部得分 0, 在 6 个真迷宫生成器上均值 0.711, 间隔 +0.711, 误判 0/9。对同一数据集 Bellot 2021 $F = \nu/\delta$ 真伪间隔 $-191.3$, 方向反转, 误判 4/9 (Spiral $F{=}11.85$ 与 3 个 $F{=}0$ 图案排名低于真迷宫最低值 $F{=}15.33$)。$maze\_quality$ 的 $M_s$ 显式惩罚单路径占比过高, $M_j$ 显式奖励分叉比例 $\approx 0.10$, 二者共同拒收 Spiral。

\noindent\textbf{两个尺度的演化扫描.}
小范围 sweep 122 次 ($200 \times 500$) 与大范围 sweep 5/5 ($500 \times 2000$) 协同揭示: manhattan-2 / maxFam=8 配置在 $40 \times 60$ 上已接近理论上界 (小范围峰值 0.8233, 大范围峰值 0.8195, 差 0.0037)。族数上限实验 ``均值降 / 峰值升'' 现象表明: maxFam=8 的均值 (0.6306) 差于 maxFam=1 (0.6984), 但拥有全场最高 0.8233 --- 多族容量升高拖累均值, 但保留了 ``偶发突破'' 的尾部峰值。

\noindent\textbf{两类失败模式.}
(1) 预算受限型 (chebyshev-4, 0.157): $2^{103}$ 搜索空间 vs $10^5$ 预算, 高维无梯度平台; (2) 表征受限型 (manhattan-1, 0.421): 4 邻域无法稳定涌现走廊结构。二者分属不同类别, 混淆会误导调参方向: 预算受限型可通过增加预算或降维解决, 表征受限型必须换模板。

\noindent\textbf{数据可用性.}
所有代码 (maze-web 浏览器端 ES + WebGPU compute shader + 15-pattern 测试集), 训练日志 (sweep ndjson, 小范围 122 OK + 6 TIMEOUT + 大范围 5/5), 截图 (figures/v2/) 均随论文开源。15-pattern 对照基准的图案 (\texttt{paper/data/\_grids.json}) 与分数 (\texttt{paper/data/sweep\_summary.json::15pattern.per\_pattern}) 完全独立可复现, 复核脚本 \texttt{paper/data/\_eval\_15patterns.mjs} 用真实源码 (mazeQuality / bellotF) 重算, 15/15 数值精确匹配 (详见附录 B)。
